\documentclass{standalone}
\usepackage{ctex}
\usepackage{tikz}

\usetikzlibrary{cartes}

% Styles
\tikzstyle{axes}=[]
\tikzstyle{important line}=[very thick]
\tikzstyle{information text}=[rounded corners,fill=red!10,inner sep=1ex]

\begin{document}

\begin{tikzpicture}
  % Pascal's Theorem
  % The theorem states that if a six-sided polygon (hexagon) is inscribed 
  % in a conic section (such as a circle, ellipse, parabola, or hyperbola), 
  % then the three pairs of opposite sides intersect at points that lie on 
  % a single straight line. This line is called the Pascal Line.
  \tikzmath{
    % define 5 points in homogenous coordinates
    \Pa{1,1} = -2; \Pa{2,1} = 4; \Pa{3,1} = 1;
    \Pb{1,1} = -1; \Pb{2,1} = -2; \Pb{3,1} = 1;
    \Pc{1,1} = 3; \Pc{2,1} = 3; \Pc{3,1} = 1;
    \Pd{1,1} = -4; \Pd{2,1} = 1; \Pd{3,1} = 1;
    \Pe{1,1} = 1; \Pe{2,1} = 5; \Pe{3,1} = 1;
  }
  \tikzset{
    % define a conic through the 5 points above
    conics/define/five points={\Pa,\Pb,\Pc,\Pd,\Pe,\C},
    % construct the 6th point on the conic
    conics/midpoint={\Pb,\Pc,\M},% midpoint of P1 and P2
    mat/normalize={\Pa,3,1,\R},% normalize to avoid 'Dimension too large' error
    mat/normalize={\M,3,1,\S},
    conics/intersect cl={\C,\R,\S,\U,\V},% intersect line RS with the conic
    conics/exclude={\U,\V,\Pa,\Pf},% select the intersection point other than P1.
    % rearrange the 6 points if necessary
    % mat/assign={\Pa,3,1,\Temp},
    % mat/assign={\Pf,3,1,\Pa},
    % mat/assign={\Temp,3,1,\Pf},
    % construct hexagon sides
    conics/cross={\Pa,\Pb,\La},%side1: 1 -- 2
    conics/cross={\Pb,\Pc,\Lb},%side2: 2 -- 3
    conics/cross={\Pc,\Pd,\Lc},%side3: 3 -- 4
    conics/cross={\Pd,\Pe,\Ld},%side4: 4 -- 5
    conics/cross={\Pe,\Pf,\Le},%side5: 5 -- 6
    conics/cross={\Pf,\Pa,\Lf},%side6: 6 -- 1
    % intersections of opposite sides
    conics/cross={\La,\Ld,\D},%side 1 and side 4
    conics/cross={\Lb,\Le,\E},%side 2 and side 5
    conics/cross={\Lc,\Lf,\F},%side 3 and side 6
    % dehomogenize
    conics/dehomogenize={\Pa,\PPa},
    conics/dehomogenize={\Pb,\PPb},
    conics/dehomogenize={\Pc,\PPc},
    conics/dehomogenize={\Pd,\PPd},
    conics/dehomogenize={\Pe,\PPe},
    conics/dehomogenize={\Pf,\PPf},
    conics/dehomogenize={\D,\DD},
    conics/dehomogenize={\E,\EE},
    conics/dehomogenize={\F,\FF},
  }

  \coordinate (P1) at (\PPa{1},\PPa{2});
  \coordinate (P2) at (\PPb{1},\PPb{2});
  \coordinate (P3) at (\PPc{1},\PPc{2});
  \coordinate (P4) at (\PPd{1},\PPd{2});
  \coordinate (P5) at (\PPe{1},\PPe{2});
  \coordinate (P6) at (\PPf{1},\PPf{2});

  \coordinate (D) at (\DD{1},\DD{2});
  \coordinate (E) at (\EE{1},\EE{2});
  \coordinate (F) at (\FF{1},\FF{2});

  % \axes[grid] (-4:4,-2:5);
  \conic[draw=teal,thick] (\C);
  \draw (P1) -- (P2) -- (P3) -- (P4) -- (P5) -- (P6) -- cycle;
  \draw[red,thick] ($(D)!-.5!(E)$) -- ($(D)!1.5!(E)$);

  \foreach \p/\placement in {P1/above left,P2/below,P3/above right,
  P4/left,P5/above,P6/below right,D/left,E/above right,F/below}{
    \fill[red] (\p) circle (1.5pt);
    \draw (\p) node[\placement] {$\p$};
  }

  \draw[xshift=-6.15cm,yshift=-4.5cm] node [right,text width=12cm,style=information text]
    {
      \textsf{{\color{red}{帕斯卡定理}} 如果一个六边形内接于一条圆锥曲线（圆、椭圆、抛物线、双曲线），那么它的三组对边的交点共线。}

      \textit{帕斯卡定理中的六边形是指按顺序将各点依次连接（不一定凸，允许自相交）。改变顶点的排列顺序，得到的对边配对不同，因此三对边的交点也会变化。}
    };
\end{tikzpicture}

\end{document}